\section{Experiment}

In this section, we evaluate whether \model can support the generation of practice assignments from learner intent.
As discussed in the previous sections, this setting differs from conventional repository generation because the output is not only expected to be executable, but also to function as a learning environment with aligned materials, scaffolding, validation, and learner-owned work.
We therefore organize our evaluation around four research questions:
\begin{list}{$\bullet$}{%
    \setlength{\leftmargin}{1em}
    \setlength{\labelwidth}{0.5em}
    \setlength{\labelsep}{0.3em}
    \setlength{\itemsep}{0.2em}
    \setlength{\topsep}{0.2em}
    \setlength{\parsep}{0pt}
}
    \item \textbf{RQ1: Repository validity.} Can \model generate complete and executable practice-assignment repositories?
    \item \textbf{RQ2: Pedagogical alignment.} Are the generated repositories aligned with the learner's intent across project materials, scaffolding, validation, and instructor guidance?
    \item \textbf{RQ3: Learner-owned work preservation.} Does \model preserve meaningful learner-owned work while providing sufficient support for the assignment to be actionable?
    \item \textbf{RQ4: Phase-gated planning.} How much does the phase-gated workflow contribute compared with direct generation?
\end{list}
